473
8.4 Real-valued Functions of Several Variables
Fig. 8.2
and the hypothesis that
$$\frac{\partial f}{\partial x}(x_0, y_0)\xi + \frac{\partial f}{\partial y}(x_0, y_0)\eta - \zeta = 0.$$
We conclude that
$$\dot{x}(0), \dot{y}(0), \dot{z}(0)) = (\xi, \eta, \zeta).$$
Hence the tangent plane to the surface $S$ at the point $(x_0, y_0, z_0)$ is formed by the
vectors that are tangents at the point $(x_0, y_0, z_0)$ to curves on the surface $S$ passing
through the point (see Fig. 8.2).
This is a more geometric description of the tangent plane. In any case, one can
see from it that if the tangent to a curve is invariantly defined (with respect to the
choice of coordinates), then the tangent plane is also invariantly defined.
We have been considering functions of two variables for the sake of visualiz-
ability, but everything that was said obviously carries over to the general case of a
function
$$y = f(x^1,\ldots,x^m) \quad (8.77)$$
of $m$ variables, where $m \in \mathbb{N}.$
At the point $(x_0^1, \ldots, x_0^m, f (x_0^1, \ldots, x_0^m))$ the plane tangent to the graph of such
a function can be written in the form
$$y = f(x_0^1,\ldots,x_0^m) + \sum_{i=1}^m \frac{\partial f}{\partial x^i}(x_0^1,\ldots,x_0^m)(x^i - x_0^i); \quad (8.78)$$
the vector
$$\left( \frac{\partial f}{\partial x^1}(x_0), \ldots, \frac{\partial f}{\partial x^m}(x_0), -1 \right)$$
is the normal vector to the plane (8.78). This plane itself, like the graph of the func-
tion (8.77), has dimension $m$, that is, any point is now given by a set $(x^1, \ldots, x^m)$
of $m$ coordinates.
Thus, Eq. (8.78) defines a hyperplane in $\mathbb{R}^{m+1}.$
Repeating verbatim the reasoning above, one can verify that the tangent plane
(8.78) consists of vectors that are tangent to curves passing through the point